\csname documentclass\endcsname[../main.tex]{subfiles}

\begin{document}
\chapter{Implementation and Integration} \label{Implementation and Integration}

\section{Implementation Strategy}
The development of the Triplo application followed a structured and incremental approach, aligned with the architectural design of the system. The implementation was carried out using a bottom-up strategy, where foundational components were developed first and progressively extended towards higher-level functionalities. \\
In particular, the implementation process evolved through a set of stages that were not strictly sequential but partially overlapped during development. \\
First, the \textbf{model components} were defined. These components represent the core domain entities of the application and include \texttt{challengeModel}, \texttt{trekkingModel}, \texttt{UserModel}, and \texttt{DiaryModel}. They provide a structured representation of the data managed by the system and exchanged with the backend. \\
Second, the \textbf{controller components} were introduced. Controllers were developed early in the process and progressively refined as new functionalities were added. Rather than being completed in a single phase, they evolved incrementally together with the introduction of application features. This includes \texttt{languagecontroller}, \texttt{challengecontroller}, \texttt{trekkingcontroller}, \texttt{UserController}, and \texttt{DiaryController}. Over time, controllers were improved to better coordinate logic and interactions among components. \\
As development progressed, the implementation of the \textbf{view components} and controllers proceeded together, following the introduction of new functionalities. The user interface components, such as \texttt{GeowatchPage}, \texttt{HomePage}, \texttt{trekkingPage}, \texttt{SettingsPage}, \texttt{LoginRegistrationPage}, \texttt{SearchPage}, \texttt{UserProfilePage}, and \texttt{DiaryPage}, were developed incrementally to support the growing feature set of the application. \\
Finally, the \textbf{service components} were introduced. Services such as \texttt{Background\_service}, \texttt{InternetService}, \texttt{permissionservice}, \texttt{geoservice}, \texttt{memoryservice}, \texttt{notificationservice}, and \texttt{authservice} were implemented to reduce the responsibilities of controllers and to delineate and isolate individual technical duties. In particular, services encapsulate interactions with the operating system, for example permissions, geolocation, notifications, and provide reusable functionality. This design improves modularity and maintainability by decoupling application logic from platform-specific operations. \\

\newpage
\subsection{Implementation Process Overview}
The incremental construction of the system components is summarized in Figure~\ref{Implementation Process}







\label{Implementation Process}
\begin{figure}[H]
    \centering
    \includegraphics[width=0.6\linewidth]{./images/Implementation process.jpg}
    \caption{Implementation Process}
\end{figure}.


\section{Integration of Software Components}

After the implementation of the individual components, the system was progressively assembled through an integration phase. This phase followed a clear incremental strategy.
Initially, \textbf{Firebase} was integrated together with the controller components. Controllers directly interact with Firebase services (such as authentication, database, and storage) through platform-specific SDKs, enabling early validation of backend-related functionalities.
Subsequently, the integration of \textbf{view components} and controllers continued in parallel with the implementation of functionalities. This allowed the application to evolve incrementally, ensuring that each feature was fully integrated before introducing new ones.
External services, such as weather and map providers, were integrated at a later stage of development, once the corresponding user interface components were available. As the number of external dependencies increased, a dedicated \textbf{ServiceController} was introduced following the Facade design pattern. This controller centralizes access to external APIs and OS-level services, simplifying interactions and reducing coupling across the system.
Finally, \textbf{device-level services} were fully integrated through the service components, which manage interactions with the operating system in a controlled and modular way.
The integration process was incremental: core backend functionality was integrated first, followed by user interface components, and finally external services and OS interactions.










\label{Package Diagram Smartphone}
\begin{figure}[H]
    \centering
    \includegraphics[width=\linewidth]{./images/package diagram.drawio.png}
    \caption{Package diagram representing the dependencies in the smartphone application}
\end{figure}

\begin{figure}[H]
    \label{Package Diagram Smartwatch}
    \centering
    \includegraphics[width=\linewidth]{./images/diagramma-Page-2.drawio.png}
    \caption{Package Diagram representing the dependencies in the smartwatch application}
\end{figure}


\newpage
\section{Testing}

\subsection{Objectives}
The purpose of the tests written for this application is to verify logical correctness of the components and their behavior, both normal behavior and error handling.
In order to test the behavior of components that use backend funcitonalities or external APIs, those functionalities have been mocked through the \textit{Mockito} library, allowing the tests to run offline without the need of a network connection: furthermore, this choice allows to test the application logic independently from the backend and the APIs. \\
The UI, the models and the controller and services funcitonalities are tested using a mixed approach, mostly by the use of unit and widget tests. 
The widget tests are useful in order to verify the correct behavior of the UI, the navigation and response to user input.
The tests of the model components verify the correct behavior of the serialization and deserialization of the objects.
The tests of the controllers verify the bahavior of the application logic without the need of the backend or APIs. 

\subsection{Application tests coverage}

\begin{figure}[H]
    \label{test smartphone}
    \centering
    \includegraphics[width=\linewidth]{./images/test smartphone.jpg}
    \caption{test smartphone application}
\end{figure}


\newpage
\begin{figure}[H]
    \label{test smartwatch}
    \centering
    \includegraphics[width=0.8\linewidth]{./images/test smartwatch.jpg}
    \caption{test smartwatch application}
\end{figure}


\subsection{Widget Tests}
The widget tests were used to check if the UI interface behaves in the correct way. In this appliaction they are very important, because many pages contain buttons, dialogs and content that changes depending on the app state.
These tests aim to verify if a page is rendered correctly and if the user is moved to the right page in response to the user's taps. 
For example, the tests for the search page on the smartphone application, among the graphical elements of the page, check the search field, the user and trekking filters. 
Other widget tests, like the ones for the settings page or the trekking page, check forms, dialogs, profile actions and screen navigation.
\subsection{Module Tests}
The module tests are used to verify larger parts of the application logic, not only a single small function. In this project, they are mainly used for controllers and services, because these components manage many operations together.
A module test usually checks if one component works correctly when it receives some input and interacts with mocked dependencies. For example, the authentication service tests of the smartphone application verify registration, login, logout, password reset and watch connection procedures, the user controller tests check if the controller updates the internal state correctly after login, profile update or follower actions.
These tests are important because they show if one application module behaves well as a whole, even without using the real backend.

\subsubsection{Model Tests}
The model tests are focused on data structures. Their goal is to verify if the application reads, stores and converts data in the right way.
For example, the tests for the model check constructors, getters, setters, default values and especially the methods \texttt{toMap()} and \texttt{fromMap()}. In many cases the tests also verify what happens if a field is missing, null or written with an old key name.
\subsubsection{Controller Tests}
The controller tests are among the most useful ones in the whole project, because the controllers contain a big part of the application logic.
For example, the trekking controller tests verify trekking loading, search, saved routes. The diary controller tests are used for adding, editing and removing diaries, including some edge cases.




\subsubsection{Service Tests}
These tests are used to validate the service components of the app, dealing with authentication, memory cache, notifications, internet connection, permissions, geolocation and background tasks.
For example, the authentication service tests verify the login and registration flows with mocked Firebase objects. The memory service tests check RAM cache, saved locale values and keys. There are also tests for internet status changes, GPS access and permission handling.


\subsection{User Testing}

User testing represents the final stage of validation, where the application is tested by end users in realistic scenarios. This type of testing focuses on usability, user experience, and overall system behavior in real-world conditions.

\subsection{Test Environment}

The testing and validation activities are performed using a combination of physical devices and emulators, as summarized in Table~\ref{tab:test-environment}.

\begin{table}[h]
\centering
\begin{tabular}{|c|c|c|c|}
\hline
\textbf{Device} & \textbf{Type} & OS & \textbf{Architecture} \\
\hline
Android Emulator (Medium Phone API 36) & Emulator & Android & x86\_64 \\
Wear OS Emulator (Small Round) & Emulator & Android Wear OS & x86\_64 \\
Samsung Galaxy A34 5G & Real Device & Android & ARM64 \\
Motorola Edge 60 &  Real Device & Android & ARM 64 \\
iPhone 15 Plus & Real Device & iOS & ARM64 \\
\hline
\end{tabular}
\caption{Testing environment devices and architectures}
\label{tab:test-environment}
\end{table}







\newpage
\paragraph{Weather Alerts Demo}

For demonstration and testing purposes, weather alert notifications can be manually sent through a POST request to the demo server. This allows developers to simulate alert conditions and verify the correct behavior of the system.

On Windows (PowerShell), the following command can be used:

{\small
\begin{verbatim}
Invoke-WebRequest "https://meteodemoserver.onrender.com/alerts" `
-Method POST `
-UseBasicParsing `
-Headers @{ "Content-Type"="application/json" } `
-Body '[
  {
    "event":"Heavy Snow",
    "severity":"Severe",
    "headline":"Snowstorm",
    "description":"Up to 50cm above 1200m"
  },
  {
    "event":"Low visibility",
    "severity":"minor",
    "headline":"Low visibility",
    "description":"Up to 50cm above 1200m"
  }
]' | Select-Object -ExpandProperty Content
\end{verbatim}
}

On Linux and macOS (bash), the equivalent request can be executed using \texttt{curl}:





{\small
\begin{verbatim}
curl -X POST "https://meteodemoserver.onrender.com/alerts" \
-H "Content-Type: application/json" \
-d '[
  {
    "event":"Heavy Snow",
    "severity":"Severe",
    "headline":"Snowstorm",
    "description":"Up to 50cm above 1200m"
  },
  {
    "event":"Low visibility",
    "severity":"minor",
    "headline":"Low visibility",
    "description":"Up to 50cm above 1200m"
  }
]'
\end{verbatim}
}
A successful request returns the following response:
\begin{verbatim}
{"status":"ok"}
\end{verbatim}
This mechanism is used exclusively for testing and demonstration purposes and does not represent the production workflow of the alert system.





\end{document}
